\section{AHSME 1986}
        \begin{enumerate}
        	\item (\textit{AHSME 1986}) $[x-(y-z)] - [(x-y) - z] =$
$\textbf{(A)}\ 2y \qquad \textbf{(B)}\ 2z \qquad \textbf{(C)}\ -2y \qquad \textbf{(D)}\ -2z \qquad \textbf{(E)}\ 0$


	\item (\textit{AHSME 1986}) If the line $L$ in the $xy$-plane has half the slope and twice the $y$-intercept of the line $y = \frac{2}{3} x + 4$, then an equation for $L$ is:


 $\textbf{(A)}\ y = \frac{1}{3} x + 8 \qquad \textbf{(B)}\ y = \frac{4}{3} x + 2 \qquad \textbf{(C)}\ y =\frac{1}{3}x+4\qquad\\  \textbf{(D)}\ y =\frac{4}{3}x+4\qquad \textbf{(E)}\ y =\frac{1}{3}x+2$


	\item (\textit{AHSME 1986}) $\triangle ABC$ has a right angle at $C$ and $\angle A = 20^\circ$. If $BD$ ($D$ in $\overline{AC}$) is the bisector of $\angle ABC$, then $\angle BDC =$
$\textbf{(A)}\ 40^\circ \qquad \textbf{(B)}\ 45^\circ \qquad \textbf{(C)}\ 50^\circ \qquad \textbf{(D)}\ 55^\circ\qquad \textbf{(E)}\ 60^\circ$


	\item (\textit{AHSME 1986}) Let S be the statement 
"If the sum of the digits of the whole number $n$ is divisible by $6$, then $n$ is divisible by $6$."


 A value of $n$ which shows $S$ to be false is


 $\textbf{(A)}\ 30 \qquad \textbf{(B)}\ 33 \qquad \textbf{(C)}\ 40 \qquad \textbf{(D)}\ 42 \qquad \textbf{(E)}\ \text{ none of these}$


	\item (\textit{AHSME 1986}) Simplify $\left(\sqrt[6]{27} - \sqrt{6 \frac{3}{4} }\right)^2$
$\textbf{(A)}\ \frac{3}{4} \qquad \textbf{(B)}\ \frac{\sqrt 3}{2} \qquad \textbf{(C)}\ \frac{3\sqrt 3}{4}\qquad \textbf{(D)}\ \frac{3}{2}\qquad \textbf{(E)}\ \frac{3\sqrt 3}{2}$


	\item (\textit{AHSME 1986}) Using a table of a certain height, two identical blocks of wood are placed as shown in Figure 1.  Length $r$ is found to be $32$ inches.  After rearranging the blocks as in Figure 2, length $s$ is found to be $28$ inches.  How high is the table?


 
                    \begin{figure}[H]
                        \centering
                        \includegraphics[width=0.5\textwidth]{Problem/AHSME_I_1986_6_Problem_Figure_1.png}
                    \end{figure}
                    
$\textbf{(A) }28\text{ inches}\qquad\textbf{(B) }29\text{ inches}\qquad\textbf{(C) }30\text{ inches}\qquad\textbf{(D) }31\text{ inches}\qquad\textbf{(E) }32\text{ inches}$


	\item (\textit{AHSME 1986}) The sum of the greatest integer less than or equal to $x$ and the least integer greater than or equal to $x$ is $5$. The solution set for $x$ is


 $\textbf{(A)}\ \Big\{\frac{5}{2}\Big\}\qquad \textbf{(B)}\ \big\{x\ |\ 2 \le x \le 3\big\}\qquad \textbf{(C)}\ \big\{x\ |\ 2\le x < 3\big\}\qquad\\  \textbf{(D)}\ \Big\{x\ |\ 2 < x\le 3\Big\}\qquad \textbf{(E)}\ \Big\{x\ |\ 2 < x < 3\Big\}$


	\item (\textit{AHSME 1986}) The population of the United States in $1980$ was $226,504,825$. The area of the country is $3,615,122$ square miles. There are $(5280)^{2}$ 
square feet in one square mile. Which number below best approximates the average number of square feet per person?


 $\textbf{(A)}\ 5,000\qquad \textbf{(B)}\ 10,000\qquad \textbf{(C)}\ 50,000\qquad \textbf{(D)}\ 100,000\qquad \textbf{(E)}\ 500,000$


	\item (\textit{AHSME 1986}) The product $\left(1-\frac{1}{2^{2}}\right)\left(1-\frac{1}{3^{2}}\right)\ldots\left(1-\frac{1}{9^{2}}\right)\left(1-\frac{1}{10^{2}}\right)$ equals


 $\textbf{(A)}\ \frac{5}{12}\qquad \textbf{(B)}\ \frac{1}{2}\qquad \textbf{(C)}\ \frac{11}{20}\qquad \textbf{(D)}\ \frac{2}{3}\qquad \textbf{(E)}\ \frac{7}{10}$


	\item (\textit{AHSME 1986}) The $120$ permutations of $AHSME$ are arranged in dictionary order as if each were an ordinary five-letter word. 
The last letter of the $86$th word in this list is:


 $\textbf{(A)}\ \text{A} \qquad \textbf{(B)}\ \text{H} \qquad \textbf{(C)}\ \text{S} \qquad \textbf{(D)}\ \text{M}\qquad \textbf{(E)}\ \text{E}$


	\item (\textit{AHSME 1986}) In $\triangle ABC, AB = 13, BC = 14$ and $CA = 15$. Also, $M$ is the midpoint of side $AB$ and $H$ is the foot of the altitude from $A$ to $BC$. 
The length of $HM$ is


 
                    \begin{figure}[H]
                        \centering
                        \includegraphics[width=0.5\textwidth]{Problem/AHSME_I_1986_11_Problem_Figure_1.png}
                    \end{figure}
                    
$\textbf{(A)}\ 6\qquad \textbf{(B)}\ 6.5\qquad \textbf{(C)}\ 7\qquad \textbf{(D)}\ 7.5\qquad \textbf{(E)}\ 8$


	\item (\textit{AHSME 1986}) John scores $93$ on this year's AHSME. Had the old scoring system still been in effect, he would score only $84$ for the same answers. 
How many questions does he leave unanswered? (In the new scoring system that year, one received $5$ points for each correct answer,
$0$ points for each wrong answer, and $2$ points for each problem left unanswered. In the previous scoring system, 
one started with $30$ points, received $4$ more for each correct answer, 
lost $1$ point for each wrong answer, and neither gained nor lost points for unanswered questions.)


 $\textbf{(A)}\ 6\qquad \textbf{(B)}\ 9\qquad \textbf{(C)}\ 11\qquad \textbf{(D)}\ 14\qquad \textbf{(E)}\ \text{Not uniquely determined}$


	\item (\textit{AHSME 1986}) A parabola $y = ax^{2} + bx + c$ has vertex $(4,2)$. If $(2,0)$ is on the parabola, then $abc$ equals


 $\textbf{(A)}\ -12\qquad \textbf{(B)}\ -6\qquad \textbf{(C)}\ 0\qquad \textbf{(D)}\ 6\qquad \textbf{(E)}\ 12$


	\item (\textit{AHSME 1986}) Suppose hops, skips and jumps are specific units of length. If $b$ hops equals $c$ skips, $d$ jumps equals $e$ hops, 
and $f$ jumps equals $g$ meters, then one meter equals how many skips?


 $\textbf{(A)}\ \frac{bdg}{cef}\qquad \textbf{(B)}\ \frac{cdf}{beg}\qquad \textbf{(C)}\ \frac{cdg}{bef}\qquad \textbf{(D)}\ \frac{cef}{bdg}\qquad \textbf{(E)}\ \frac{ceg}{bdf}$


	\item (\textit{AHSME 1986}) A student attempted to compute the average, $A$, of $x, y$ and $z$ by computing the average of $x$ and $y$, and then computing the average of the result and $z$. Whenever $x < y < z$, the student's final result is


 $\textbf{(A)}\ \text{correct}\quad \textbf{(B)}\ \text{always less than A}\quad \textbf{(C)}\ \text{always greater than A}\quad\\ \textbf{(D)}\ \text{sometimes less than A and sometimes equal to A}\quad\\ \textbf{(E)}\ \text{sometimes greater than A and sometimes equal to A}  \quad$


	\item (\textit{AHSME 1986}) In $\triangle ABC, AB = 8, BC = 7, CA = 6$ and side $BC$ is extended, as shown in the figure, to a point $P$ so that $\triangle PAB$
is similar to $\triangle PCA$. The length of $PC$ is


 
                    \begin{figure}[H]
                        \centering
                        \includegraphics[width=0.5\textwidth]{Problem/AHSME_I_1986_16_Problem_Figure_1.png}
                    \end{figure}
                    
$\textbf{(A)}\ 7\qquad \textbf{(B)}\ 8\qquad \textbf{(C)}\ 9\qquad \textbf{(D)}\ 10\qquad \textbf{(E)}\ 11$


	\item (\textit{AHSME 1986}) A drawer in a darkened room contains $100$ red socks, $80$ green socks, $60$ blue socks and $40$ black socks. 
A youngster selects socks one at a time from the drawer but is unable to see the color of the socks drawn. 
What is the smallest number of socks that must be selected to guarantee that the selection contains at least $10$ pairs? 
(A pair of socks is two socks of the same color. No sock may be counted in more than one pair.)


 $\textbf{(A)}\ 21\qquad \textbf{(B)}\ 23\qquad \textbf{(C)}\ 24\qquad \textbf{(D)}\ 30\qquad \textbf{(E)}\ 50$


	\item (\textit{AHSME 1986}) A plane intersects a right circular cylinder of radius $1$ forming an ellipse. 
If the major axis of the ellipse is $50\%$ longer than the minor axis, the length of the major axis is


 $\textbf{(A)}\ 1\qquad \textbf{(B)}\ \frac{3}{2}\qquad \textbf{(C)}\ 2\qquad \textbf{(D)}\ \frac{9}{4}\qquad \textbf{(E)}\ 3$


	\item (\textit{AHSME 1986}) A park is in the shape of a regular hexagon $2$ km on a side. Starting at a corner, 
Alice walks along the perimeter of the park for a distance of $5$ km. 
How many kilometers is she from her starting point?


 $\textbf{(A)}\ \sqrt{13}\qquad \textbf{(B)}\ \sqrt{14}\qquad \textbf{(C)}\ \sqrt{15}\qquad \textbf{(D)}\ \sqrt{16}\qquad \textbf{(E)}\ \sqrt{17}$


	\item (\textit{AHSME 1986}) Suppose $x$ and $y$ are inversely proportional and positive. If $x$ increases by $p\%$, then $y$ decreases by 


 $\textbf{(A)}\ p\%\qquad \textbf{(B)}\ \frac{p}{1+p}\%\qquad \textbf{(C)}\ \frac{100}{p}\%\qquad \textbf{(D)}\ \frac{p}{100+p}\%\qquad \textbf{(E)}\ \frac{100p}{100+p}\%$


	\item (\textit{AHSME 1986}) In the configuration below, $\theta$ is measured in radians, $C$ is the center of the circle, 
$BCD$ and $ACE$ are line segments and $AB$ is tangent to the circle at $A$.


 
                    \begin{figure}[H]
                        \centering
                        \includegraphics[width=0.5\textwidth]{Problem/AHSME_I_1986_21_Problem_Figure_1.png}
                    \end{figure}
                    


 A necessary and sufficient condition for the equality of the two shaded areas, given $0 < \theta < \frac{\pi}{2}$, is


 $\textbf{(A)}\ \tan \theta = \theta\qquad \textbf{(B)}\ \tan \theta = 2\theta\qquad \textbf{(C)}\ \tan\theta = 4\theta\qquad \textbf{(D)}\ \tan 2\theta =\theta\qquad\\  \textbf{(E)}\ \tan\frac{\theta}{2}=\theta$


	\item (\textit{AHSME 1986}) Six distinct integers are picked at random from $\{1,2,3,\ldots,10\}$. What is the probability that, among those selected, the second smallest is $3$?


 $\textbf{(A)}\ \frac{1}{60}\qquad \textbf{(B)}\ \frac{1}{6}\qquad \textbf{(C)}\ \frac{1}{3}\qquad \textbf{(D)}\ \frac{1}{2}\qquad \textbf{(E)}\ \text{none of these}$


	\item (\textit{AHSME 1986}) Let $N = 69^{5} + 5 \cdot 69^{4} + 10 \cdot 69^{3} + 10 \cdot 69^{2} + 5 \cdot 69 + 1$. How many positive integers are factors of $N$?


 $\textbf{(A)}\ 3\qquad \textbf{(B)}\ 5\qquad \textbf{(C)}\ 69\qquad \textbf{(D)}\ 125\qquad \textbf{(E)}\ 216$


	\item (\textit{AHSME 1986}) Let $p(x) = x^{2} + bx + c$, where $b$ and $c$ are integers. 
If $p(x)$ is a factor of both $x^{4} + 6x^{2} + 25$ and $3x^{4} + 4x^{2} + 28x + 5$, what is $p(1)$?


 $\textbf{(A)}\ 0\qquad \textbf{(B)}\ 1\qquad \textbf{(C)}\ 2\qquad \textbf{(D)}\ 4\qquad \textbf{(E)}\ 8$


	\item (\textit{AHSME 1986}) If $\lfloor x\rfloor$ is the greatest integer less than or equal to $x$, then
$\sum_{N=1}^{1024} \lfloor \log_{2}N\rfloor =$
$\textbf{(A)}\ 8192\qquad \textbf{(B)}\ 8204\qquad \textbf{(C)}\ 9218\qquad \textbf{(D)}\ \lfloor\log_{2}(1024!)\rfloor\qquad \textbf{(E)}\ \text{none of these}$


	\item (\textit{AHSME 1986}) It is desired to construct a right triangle in the coordinate plane so that its legs are parallel 
to the $x$ and $y$ axes and so that the medians to the midpoints of the legs lie on the lines $y = 3x + 1$
and $y = mx + 2$. The number of different constants $m$ for which such a triangle exists is


 $\textbf{(A)}\ 0\qquad \textbf{(B)}\ 1\qquad \textbf{(C)}\ 2\qquad \textbf{(D)}\ 3\qquad \textbf{(E)}\ \text{more than 3}$


	\item (\textit{AHSME 1986}) In the adjoining figure, $AB$ is a diameter of the circle, $CD$ is a chord parallel to $AB$, 
and $AC$ intersects $BD$ at $E$, with $\angle AED = \alpha$. The ratio of the area of $\triangle CDE$ to that of $\triangle ABE$ is  


 
                    \begin{figure}[H]
                        \centering
                        \includegraphics[width=0.5\textwidth]{Problem/AHSME_I_1986_27_Problem_Figure_1.png}
                    \end{figure}
                    
$\textbf{(A)}\ \cos\ \alpha\qquad \textbf{(B)}\ \sin\ \alpha\qquad \textbf{(C)}\ \cos^2\alpha\qquad \textbf{(D)}\ \sin^2\alpha\qquad \textbf{(E)}\ 1-\sin\ \alpha$


	\item (\textit{AHSME 1986}) $ABCDE$ is a regular pentagon. $AP, AQ$ and $AR$ are the perpendiculars dropped from $A$ onto $CD, CB$ extended and $DE$ extended, 
respectively. Let $O$ be the center of the pentagon. If $OP = 1$, then $AO + AQ + AR$ equals


 
                    \begin{figure}[H]
                        \centering
                        \includegraphics[width=0.5\textwidth]{Problem/AHSME_I_1986_28_Problem_Figure_1.png}
                    \end{figure}
                    
$\textbf{(A)}\ 3\qquad \textbf{(B)}\ 1 + \sqrt{5}\qquad \textbf{(C)}\ 4\qquad \textbf{(D)}\ 2 + \sqrt{5}\qquad \textbf{(E)}\ 5$


	\item (\textit{AHSME 1986}) Two of the altitudes of the scalene triangle $ABC$ have length $4$ and $12$. 
If the length of the third altitude is also an integer, what is the biggest it can be?


 $\textbf{(A)}\ 4\qquad \textbf{(B)}\ 5\qquad \textbf{(C)}\ 6\qquad \textbf{(D)}\ 7\qquad \textbf{(E)}\ \text{none of these}$


	\item (\textit{AHSME 1986}) The number of real solutions $(x,y,z,w)$ of the simultaneous equations
$2y = x + \frac{17}{x}, 2z = y + \frac{17}{y}, 2w = z + \frac{17}{z}, 2x = w + \frac{17}{w}$
is


 $\textbf{(A)}\ 1\qquad \textbf{(B)}\ 2\qquad \textbf{(C)}\ 4\qquad \textbf{(D)}\ 8\qquad \textbf{(E)}\ 16$


\end{enumerate}